\scanpage{92}


i.e.

\[ 4c=4\left(b_1^2c_2+b_2^2c_1-4c_1c_2\right)=4(c_1\delta_2+c_2\delta_1+4c_1c_2) \]

\inserted{(\uline{\unsure{heuristiquement}})}

et la formule pertinente pour $c$ s'en déduit ``en divisant par 4''\ldots

Pour prouver le \unsure{th.|th}, il suffit de le faire localement, à cause de l'unicité, donc on peut choisir un scindage, et définir $Q$ \unsure{ad hoc|obtenue} par $b$ et $c$ donnés par les formules plus haut. Il faut prouver que $Q(y_1*y_2)$ satisfait la formule écrite -- c'est un petit calcul \unsure{délectable|détestable}, que j'ai pratiquement déjà fait plus haut.

\uline{Définition} de $X_1\unsure{\wedge|*|\star}X_2$ et de

\[ X_1\times X_2\longrightarrow X_1\wedge X_2 \]

Contraintes de comm, assoc, unité pour cette opération

$X_1\wedge X_2$ quand $X_1$ (ou $X_2$) est étale, comme déduit de $X_2$ en le ``tordant'' par le $(\unsure{\mathbb{Z}_{2\mathbb Z}|\mathbb Z/2\mathbb Z})_X$-torseur $X_1$

$X_1\wedge X_2$ est étale ssi $X_1$ et $X_2$ le sont

$X_1\times X_2\unsure{-|\to}X_1\wedge X_2$ est fid. plat, sauf
% archive sheet number: 91
